\chapter{The Big Questions and a Few Smaller Ones}
\label{chap:faq}

\subsection*{Can a mechanical system, such as a computer, feel pain?}

A computer—or even a simple thermostat—given sufficient configuration space, could in theory duplicate the exact informational state of a suffering human.
However, this does not make the physical hardware device itself conscious.
There is no room for qualia to causally affect the deterministic or statistical operation of the physical system.


\subsection*{Can a computer simulate pain?}

We, as observers, are static, finite sets of information.
Based on observational evidence, we cannot create new fundamental information—we merely explore the combinations latent within our set. 

Correspondingly, a digital simulations won't create new instance of suffering.
The simulation merely traces out one of the astronomically many configurations that already describe the same equivalence class of that observer.
Alice exists eternally within the static totality ($U$) regardless of whether a human-built machine ever executes her corresponding bitstring.


\subsection*{What explains the power of quantum computers?}

Some physicists, like David Deutsch \cite{deutsch1997fabric}, argue that quantum computers achieve their speed by performing calculations across many parallel universes simultaneously.
Under the lens of the present theory, quantum computers do not create answers through dynamic, parallel computational labor.
Instead, they act as informational filters that reveal structure already inherently present.


\subsection*{Is this a ``One-Parameter'' Theory rather than Zero?}

Strictly speaking, $\lambda$ appears as an explicit parameter in our probability measure: 
\[
\mathbb{P}(\gamma \mid O) \propto \exp\bigl[-\lambda\, \mathcal{C}_O(\gamma)\bigr]
\]
However, $\lambda$ is not an arbitrary, fine-tuned input like the mass of a Higgs boson or the cosmological constant.
It represents the Global Compression Pressure—the structural degree to which the informational substrate is packed into predictable, continuous patterns.

\paragraph{The Analogy of Surface Tension}
Think of $\lambda$ like the surface tension of a water droplet.
One does not need to manually input surface tension as an independent rule to create a drop; the tension is an emergent property of how the underlying molecules minimize their energy configurations. 

In this theory:
\begin{itemize}
    \item $\lambda$ behaves as a Selection Pressure: It dictates how strictly the universe favors compressible simplicity over high-entropy chaos.
    \item If $\lambda \to 0$, the measure flattens, and the universe dissolves into unconstrained white noise (infinite spectral complexity).
    \item If $\lambda \to \infty$, the universe collapses into a single, maximally simple static state.
\end{itemize}


\subsection*{Where did all the matter in the universe come from?}

This question implicitly assumes a fundamental temporal ordering, a causal transition, and an external point of origin.
Because both time and causality are entirely emergent, observer-dependent phenomena (as established in Chapter~\nameref{chap:humans-as-axiomatic-systems}), the question dissolves.
Matter does not ``come from'' anywhere; it coexists statically as part of the valid informational geometries that define our observer block.


\subsection*{Why is there something rather than nothing?}

This is analogous to asking why a flipped coin shows ``heads'' rather than ``tails'': both are abstract states within a configuration space, and neither is ontologically privileged.
If there were a cosmic rule that favored absolute nothingness over something, we would be forced to ask about the origin of that rule—who enforced the constraint of emptiness? Non-empty structures are just as real or unreal as the empty set.
Both exist simply because there is no external meta-rule to forbid them.


\subsection*{Why do particles follow a complex-valued wavefunction?}

Because we are observing relational structure through the lens of maximal compression.
The wavefunction is the \textbf{spectral compression codec} of reality's underlying informational structure.
Observed wave-like behavior emerges naturally when representing dense combinatorial constraints with minimal description length.

Formally, the hierarchy progresses as:
\[
\boxed{\text{Maximal Predictability}} \;\rightarrow\; \boxed{\text{Maximal Compression}} \;\rightarrow\; \boxed{\text{Maximal Probability}} 
\]
Predictability is exactly what we perceive as the stable ``laws of physics.''
Configurations with minimal observer-indexed spectral complexity $\mathcal{C}_O[\gamma]$ overwhelmingly dominate the probability measure:
\[
\mathbb{P}(\gamma \mid O) \propto \exp\bigl[-\lambda \, \mathcal{C}_O[\gamma]\bigr].
\]
Hence, the laws of physics are not rigid commandments; they are the statistical signatures of highly compressible configurations.


\subsection*{Why does the universe expand? Is there a link to entropy?}

Spacetime expansion is the geometric manifestation of increasing informational configuration entropy.
A bitstring of zero entropy corresponds geometrically to a singularity—a mapping with zero volume.
As the entropy of the configuration increases, the corresponding geometric representation must necessarily stretch to accommodate the growing density of internal relational constraints.
Volume is merely an expression of informational capacity.


\subsection*{Why did the universe start with a low-entropy state?}

Low-entropy configurations are maximally compressible and therefore receive the highest statistical weight.
A high-entropy initial state would require an immense amount of additional information just to encode the localized observer, which drastically increases the spectral complexity $\mathcal{C}_O[\gamma]$ and suppresses its probability measure:
\[
\mathbb{P}(\gamma_{\rm high-entropy} \mid O) \ll \mathbb{P}(\gamma_{\rm low-entropy} \mid O).
\]
Consequently, the thermodynamic arrow of time emerges naturally from an origin of absolute simplicity:
\[
\text{Minimal Complexity} \;\Rightarrow\; \text{Maximal Compressibility} \;\Rightarrow\; \text{Highly Probable Observer Path}.
\]


\subsection*{Is the Universe Infinite?}

If reality were finite, one would have to ask: finite by exactly how many bits? And what external mechanism enforced that specific boundary? Any fixed numerical constraint requires an independent explanation, which merely pushes the mystery back one level. 

The only non-arbitrary conclusion is that the foundational substrate of reality cannot be finitely bounded.
At the deepest level, the static totality is infinite, existing far beyond any localized, finite description.


\subsection*{What is gravity?}

We observe gravity for the same structural reason we observe wave mechanics: we are experiencing a highly compressed representation of information. 

Observers inevitably find themselves within the most compressible (and thus most probable) global configurations.
The smooth trajectory through these configurations is what we perceive as continuous space and time.
We ``fall'' because there is more copies of us ``down there''.
Gravity is the statistical gradient driving the observer toward the most highly probable, compressible arrangements.


\subsection*{How does Gödel's incompleteness theorem affect the theory?}

Stephen Hawking (in ``Gödel and the End of Physics'') argued that because a physical Theory of Everything must describe itself, it represents a self-referential formal system and must therefore be incomplete. 

By placing the static totality $U$ entirely outside the domain of formal axiomatic systems, this framework dissolves the contradiction.
Formal math, symbolic logic, and physical equations are not the fundamental bedrock of reality; they are emergent,
interpretive compression artifacts within observer equivalence classes. Because $U$ is an uninterpreted totality rather than a formal mathematical system, it is completely immune to Gödelian incompleteness.


\subsection*{Are we living in a simulation?}

If we were, we would be forced to ask: who simulates the simulator? This framework eliminates the infinite regress.
We are simultaneously the simulation and the simulator.
The observer and the observed environment are entangled components of the exact same static informational block.


\subsection*{Does the theory support QBism?}

Quantum Bayesianism (QBism) interprets the wavefunction as a representation of an agent’s subjective degrees of belief,
viewing observation as an active, participatory event that alters the world.

The present theory rejects this agency. Here, the wavefunction possesses no independent physical reality; it is entirely emergent.
The universe is fundamentally static—there are no temporal transitions, no physical signals, and no dynamic wave collapses.
The observer does not ``alter'' reality by looking.


\subsection*{Why do we age and die?}

The observer’s continuation is limited by the spectral complexity of their state.
Let $\Sigma[\gamma]$ denote the number of independent frequency-phase components required to encode a future history $\gamma$.
The probability of an observer-compatible path scales inversely with its spectral inflation:
\[
\mathbb{P}(\gamma \mid O) \propto \exp\bigl[-\alpha \, \Sigma[\gamma]\bigr].
\]
The accumulation of internal memories, physical wear, and chaotic environmental interactions steadily increases $\Sigma[\gamma]$.
This rapidly reduces the available combinatorial measure of valid continuations.
Aging and death are not biological failures; they are the geometric and statistical manifestations of combinatorial exhaustion within a finite state space.


\subsection*{Is this just Solipsism?}

While this theory takes the observer's state as its axiomatic starting point, it does not collapse into solipsism.
The existence of multiple distinct observers is a straightforward probabilistic variable.
It is statistically impossible to argue that out of billions of valid, highly compressible human configurations, only a single isolated observer is allowed.
The combinatorial measure naturally favors a rich, multi-agent structural ecosystem over an isolated, high-complexity solipsistic exception.


\subsection*{What motivates the focus on Spectral Complexity?}

Direct observational evidence. When we peer into the micro-cosmos, reality does not present itself as discrete, isolated digital bits;
it manifests fundamentally as continuous waves.
Spectral complexity provides the precise, continuous, and computable metric needed to bridge discrete information with the wave mechanics of actual physics.


\subsection*{Does the theory explain why the cosmic expansion rate is near-critical?}

Conditional on typical observer emergence within an entropy-increasing configuration, we are overwhelmingly likely to find ourselves situated near the peak of the microstructure count (where compression efficiency is maximized).
In spacetime geometry, this exact statistical sweet spot maps directly to a near-critical expansion density ($\Omega \approx 1$).
